\documentclass[11pt]{article}
\usepackage[T1]{fontenc}
\usepackage{lmodern}
\usepackage{amsmath,amssymb,amsthm,mathtools}
\usepackage[margin=1in]{geometry}
\usepackage{microtype}
\usepackage[colorlinks=true,linkcolor=blue,citecolor=blue,urlcolor=blue]{hyperref}
\newtheorem{theorem}{Theorem}[section]
\newtheorem{lemma}[theorem]{Lemma}
\newtheorem{proposition}[theorem]{Proposition}
\theoremstyle{remark}
\newtheorem{remark}[theorem]{Remark}
\newcommand{\R}{\mathbb R}
\newcommand{\Mplus}{\mathcal M_+}
\newcommand{\norm}[1]{\left\lVert #1\right\rVert}
\newcommand{\abs}[1]{\left\lvert #1\right\rvert}
\DeclareMathOperator{\spanop}{span}
\title{Uniqueness of nonnegative heat controls at the minimum time}
\author{Kaiwen Jin\thanks{Contact: \href{mailto:yui.ui.siki@gmail.com}{yui.ui.siki@gmail.com}.}}
\date{October 2026}
\begin{document}
\maketitle
\begin{abstract}
We consider the heat equation on an interval with two nonnegative Dirichlet
boundary controls. For every pair of distinct positive constant states, we
prove that the nonnegative finite Radon measure pair attaining the minimum
time is unique, under the terminal sine-moment formulation. The two boundary
measures are equal. This answers the continuum uniqueness question posed by
Loh\'eac, Tr\'elat and Zuazua. The proof uses finite positive moment cones at
the critical time. Their extremal supporting functions have vanishing cost
on every admissible measure. Exponential interpolation changes their
exponents without increasing their magnitude, while a determinant estimate
keeps an inserted coefficient away from zero. These tests determine both the
average control and the difference between the boundary controls, including
endpoint atoms.
\end{abstract}

\section{Introduction}

Nonnegative boundary temperatures constrain how quickly a rod can be heated
or cooled. We study the heat equation
\begin{equation}\label{eq:heat}
 \begin{cases}
 y_t=y_{xx},&t>0,\quad 0<x<1,\\
 y(0,x)=a,\\
 y(t,0)=u_0(t),\qquad y(t,1)=u_1(t),
 \end{cases}
\end{equation}
where $a,b>0$, $a\ne b$, and the target is $y(H,x)=b$.
There is no upper bound on the permitted nonnegative boundary temperatures,
but each classical control is bounded on its time interval. Write
\begin{equation}\label{eq:mintime}
 T_*(a,b)=\inf\{H>0:\text{bounded }u_0,u_1\ge0
                   \text{ take }a\text{ to }b\text{ at time }H\}.
\end{equation}
Without the nonnegativity constraint, bounded null controls exist on every
positive duration. Positivity changes the minimum-time question:
Loh\'eac, Tr\'elat and Zuazua proved a finite, strictly positive waiting time
and attainment by nonnegative measure controls
\cite[Theorem 2.1]{LTZ2017}.

The remaining question is whether the fastest measure control is uniquely
determined. It was posed in \cite[Section 6]{LTZ2017}, with the terminal
state expressed through sine moments. Further continuum structure was
obtained by Loh\'eac for a model with one controlled
endpoint and a homogeneous condition at the other \cite[Theorem 1]{L2021Heat}.
For reachable steady-state targets generated by a positive boundary control,
that work constructs a countably impulsive minimum-time control and proves
uniqueness within the purely impulsive class. Uniqueness among all Radon
controls is left open there \cite[Section 5]{L2021Heat}.
For finite-dimensional systems with real spectral hypotheses, uniqueness was proved in
\cite{LTZ2021}. They include discretized heat systems, but finite-dimensional
uniqueness does not identify the limiting pair for the continuum equation.
Here we prove uniqueness for the full two-boundary measure pair directly from
all its moments.

We use the closed-interval measure formulation to state precisely what
attainment means. A pair $(\mu_0,\mu_1)$ of finite nonnegative Radon measures
on $[0,H]$ is admissible if, for every integer $k\ge1$,
\begin{equation}\label{eq:target}
 \int_{[0,H]}e^{-k^2\pi^2(H-t)}
       \bigl(d\mu_0(t)-(-1)^k d\mu_1(t)\bigr)
 =\frac{(1-(-1)^k)(b-ae^{-k^2\pi^2H})}{k^2\pi^2}.
\end{equation}
These are the terminal-state sine equations of \eqref{eq:heat}. For bounded
controls the measures are $u_i(t)\,dt$; the definition also retains possible
atoms at zero and at $H$. It is the explicit moment convention in AIM
Problem 138\footnote{\href{https://github.com/MColbrook/AIM/blob/c9929805cece044705abaaebe1cebddac1ba11a6/problems/138-unique-minimum-time-positive-heat-control.md}{AIM Problem 138, pinned statement.}}, so the theorem addresses that measure question without adding
a terminal weak-solution convention.

\begin{theorem}\label{thm:main}
For every $a,b>0$ with $a\ne b$, $0<T_*(a,b)<\infty$. At
$H=T_*(a,b)$ there exists exactly one nonnegative finite Radon pair on
$[0,H]$ satisfying \eqref{eq:target} for all $k\ge1$. Its components
coincide: $\mu_0=\mu_1$.
\end{theorem}

The existence and positive waiting time are inherited from the established
control theory after the convention bridge proved below. The new conclusion
is uniqueness of the complete pair. In particular, symmetry of the equation
is forced on every admissible minimum-time pair, rather than imposed on the
controls at the outset.

Two issues obstruct a direct moment argument. The average boundary measure
is prescribed only at the odd heat eigenvalues, whose reciprocals have a
finite sum; ordinary density of all powers cannot be used to determine it.
Even uniqueness of this average would leave open whether an asymmetric pair
has a signed difference with all even moments zero. Criticality supplies the
additional information needed for both questions. A positive interval of
density would allow the controls to be advanced slightly in time and corrected
by a small bounded null control, contradicting minimality.

At the fixed critical time, the finite moment cones therefore approach an
extremal initial-to-target ratio. Their nonnegative supporting exponential
polynomials have integrals tending to zero against every critical measure.
Their zeros exhaust the finite exponential zero bound. This permits us to
insert a whole interval of new exponents while retaining pointwise domination
by the original tests. Compactness results for sparse power polynomials
\cite{BE1997}, together with a normalized determinant estimate, ensure that
the new coefficient does not vanish. The resulting transforms determine the
average measure. Interlacing of odd and even eigenvalues lets the same
interpolation determine the boundary difference.

The argument concerns distinct positive constant states on the interval.
It gives neither a formula for the minimum time nor a uniqueness theorem for
nonconstant profiles. Section~\ref{sec:critical} establishes the critical
moment formulation and the no-density property. Section~\ref{sec:interpolation}
proves the interpolation estimates. Section~\ref{sec:uniqueness} constructs
the extremal tests and proves the two rigidity statements. All measure
integrals over a displayed time interval include its endpoints.

\section{Critical heat moments}\label{sec:critical}

We write $\Mplus([0,H])$ for the finite nonnegative Radon measures on the
closed interval $[0,H]$. All signed measures below are real. After dividing
the state and controls by $b$, put
\begin{equation}\label{eq:rates}
 r=\frac ab,\qquad p=\pi^2,\qquad
 \lambda_j=(2j-1)^2p,\qquad \nu_j=(2j)^2p\quad(j\ge1).
\end{equation}
This scaling is invertible and preserves every feasible time. For a boundary
pair define its average and difference by
\begin{equation}\label{eq:parity}
 w=\frac{\mu_0+\mu_1}{2b},\qquad h=\frac{\mu_0-\mu_1}{2b}.
\end{equation}
The pair is nonnegative precisely when $w\ge0$ and $|h|\le w$, where $|h|$
denotes total variation. Indeed, every Borel partition of a set $A$ gives
$\sum_\ell|h(A_\ell)|\le w(A)$; taking the supremum proves the domination.
Conversely $w\pm h\ge0$ under that domination. The inverse pair is
$\mu_0=b(w+h)$, $\mu_1=b(w-h)$.

Multiplying the terminal sine equations by the terminal exponential yields
\begin{align}
 \int_{[0,H]}e^{\lambda_jt}\,dw(t)
   &=\frac{e^{\lambda_jH}-r}{\lambda_j},\label{eq:odd}\\
 \int_{[0,H]}e^{\nu_jt}\,dh(t)&=0.\label{eq:even}
\end{align}
Thus uniqueness has two parts: uniqueness of the positive odd-moment
representation, and vanishing of the dominated signed even-moment difference.

For reference, the unnormalized sine coefficient
$y_k(t)=\int_0^1y(t,x)\sin(k\pi x)\,dx$ satisfies
\begin{equation}\label{eq:sine}
 y_k(H)=e^{-k^2pH}y_k(0)+k\pi
 \int_0^H e^{-k^2p(H-t)}\bigl(u_0(t)-(-1)^ku_1(t)\bigr)\,dt.
\end{equation}
We use the sine convention and the bounded null-controllability result of
Loh\'eac--Tr\'elat--Zuazua \cite[Section 2.2 and Lemma 2.1]{LTZ2017}:
for every $L>0$ and $z\in L^2(0,1)$, real $C^1$ boundary controls take $z$
to zero in time $L$, with
\begin{equation}\label{eq:nullbound}
 \max_{i=0,1}\norm{v_i}_{\infty}\le\kappa(L)\norm{z}_2.
\end{equation}
Their Theorem 2.1 also gives a finite positive infimum for nonnegative
$L^1$ controls between distinct positive constants. The next lemma relates
that convention to bounded controls and to the closed-interval moments.

\begin{lemma}\label{lem:attainment}
The infima over bounded nonnegative controls, nonnegative $L^1$ controls,
and nonnegative measure pairs satisfying the terminal moments agree.
Their common value $T=T_*(a,b)$ lies in $(0,\infty)$, and the measure
moments are attained on $[0,T]$.
\end{lemma}
\begin{proof}
Suppose a nonnegative pair satisfies the moments at $H>0$. Extend each
measure to $\R$ by density $a$ on $(-\infty,0)$ and density $b$ on
$(H,\infty)$, retaining its entire measure on $[0,H]$. Choose a nonnegative
smooth kernel $\rho$ with integral one and compact support in $(0,\epsilon)$.
Convolution gives smooth bounded nonnegative controls on $[0,H+\epsilon]$.

For an odd rate $\lambda$ let $\widetilde u$ be the extended average.
For every $s$ in the kernel support, its three contributions give
\begin{align*}
 e^{\lambda s}\int_{[-s,H+\epsilon-s]}e^{\lambda t}\,d\widetilde u(t)
 &=\frac{e^{\lambda s}}{\lambda}
 \bigl[a(1-e^{-\lambda s})+(be^{\lambda H}-a)
       +b(e^{\lambda(H+\epsilon-s)}-e^{\lambda H})\bigr]\\
 &=\frac{be^{\lambda(H+\epsilon)}-a}{\lambda}.
\end{align*}
Integrating against $\rho(s)\,ds$ proves the smoothed odd moments by
Tonelli's theorem. The extended difference is zero off $[0,H]$, so each
smoothed even moment is the original zero moment multiplied by
$\int e^{\nu s}\rho(s)\,ds$. Fubini applies against finite total variation
on the relevant compact interval. Both original endpoint atoms are included.

These smooth controls equal $a$ near time zero and $b$ near time
$H+\epsilon$. To see that their moments identify an actual terminal state,
subtract the affine lifting
$F(t,x)=(1-x)u_0(t)+xu_1(t)$. The remainder has homogeneous boundary values
and forcing $-F_t$. Its sine series defines an $L^2$ solution: partial sums
obey
\[
 \norm{z(t)}_2\le\norm{z(0)}_2+\int_0^t\norm{F_s(s)}_2\,ds,
\]
and the same bound for tails, followed by $L^2$ completeness, proves
convergence. Formula \eqref{eq:sine} and completeness of the sine basis
then give the constant terminal state $b$. Thus every measure-feasible time
can be approximated from above by bounded-control times. Since bounded
controls are $L^1$ controls and $L^1$ controls define measures with the same
moments, all three infima agree. The cited positive finite waiting-time
theorem applies to $a\ne b$.

Choose bounded admissible pairs at times $T_n\downarrow T$, with
$T_n\le T+1$, and extend them by zero to $[0,T+1]$. Their averages $u_n$
satisfy
\[
 u_n([0,T+1])\le\int e^{pt}\,du_n(t)
   =\frac{be^{pT_n}-a}{p}.
\]
Each component is bounded by $2u_n$, so the two positive measures have
uniformly bounded mass. Weak-star sequential compactness on the common
compact interval gives a limiting pair. Positivity passes through
nonnegative continuous tests. Any compact subset of $(T,T+1]$ eventually
lies beyond the supports, and hence has zero limiting mass; the limit is
supported on $[0,T]$. Each fixed exponential is continuous, so all terminal
moments pass to the limit. This proves attainment.
\end{proof}

From now on $T=T_*(a,b)>0$. A \emph{critical scalar measure} means a
$w\in\Mplus([0,T])$ satisfying every equation \eqref{eq:odd} at $H=T$.
Such measures exist by Lemma~\ref{lem:attainment}.

\begin{lemma}\label{lem:slack}
A critical scalar measure cannot dominate $\delta\,dt$, $\delta>0$,
on an interval $(\alpha,\beta)$ with $0<\alpha<\beta<T$.
\end{lemma}
\begin{proof}
Suppose it does, and retain $m_0=w(\{0\})$. For small $\epsilon>0$ put
$H=T-\epsilon$ and translate the part on $[\epsilon,T]$ to the left:
\[
 w^\epsilon=m_0\delta_0+
      (t\mapsto t-\epsilon)_\#(w|_{[\epsilon,T]}).
\]
The required odd moment at $H$ minus its present value is
\begin{equation}\label{eq:deficit}
 C_j^\epsilon=(e^{-\lambda_j\epsilon}-1)
                  (r/\lambda_j+m_0)
       +e^{-\lambda_j\epsilon}
          \int_{(0,\epsilon)}e^{\lambda_jt}\,dw(t).
\end{equation}
The discarded interval is open; the atom at $\epsilon$ is translated and
the original atom at zero is kept. For each fixed $j$, $C_j^\epsilon\to0$,
and for all $j,\epsilon$ its magnitude is at most
$r/p+m_0+w([0,T])$.

Choose $0<A<B<T$ strictly inside $(\alpha,\beta)$. For sufficiently small
$\epsilon$, $B<H$ and $w^\epsilon\ge\delta\,dt$ on $[A,B]$. Let
$z_\epsilon\in L^2(0,1)$ have unnormalized sine coefficients
\[
 \int_0^1z_\epsilon(x)\sin(k\pi x)\,dx
 =\begin{cases}
 -2k\pi C_j^\epsilon e^{-k^2pA},&k=2j-1,\\
 0,&k\text{ even}.
 \end{cases}
\]
This is an actual real $L^2$ state, defined by the convergent orthogonal
series. Parseval's identity and dominated convergence give
\begin{equation}\label{eq:smallstate}
 \norm{z_\epsilon}_2^2
 =8\pi^2\sum_{j\ge1}(2j-1)^2|C_j^\epsilon|^2e^{-2\lambda_jA}
 \longrightarrow0.
\end{equation}
Only odd sine modes occur, so this state is invariant under $x\mapsto1-x$.
Apply \eqref{eq:nullbound} on the fixed duration $B-A$. Averaging the two
null controls leaves the odd forcing unchanged and makes the even forcing
zero. The resulting common real control $v_\epsilon$ on $[A,B]$ has
supremum norm tending to zero. Multiplying its zero terminal odd coefficient
in \eqref{eq:sine} by $e^{\lambda_jB}/(2k\pi)$ gives
\[
 0=-C_j^\epsilon+\int_A^B e^{\lambda_jt}v_\epsilon(t)\,dt.
\]
Extend $v_\epsilon$ by zero outside $[A,B]$. Once
$\norm{v_\epsilon}_\infty\le\delta$, the measure
$w^\epsilon+v_\epsilon\,dt$ is nonnegative and satisfies all the odd
moments at $H$. Use it for both normalized boundary measures; the even
moments then vanish. The smoothing in Lemma~\ref{lem:attainment}, with extra
duration $\epsilon/2$, gives bounded controls at $T-\epsilon/2$. This
contradicts the definition of $T$.
\end{proof}

\section{Exponential interpolation}\label{sec:interpolation}

We need to change a finite test function's exponents while keeping its zeros
and controlling its magnitude. Zeros are counted with multiplicity throughout
this section.

\begin{lemma}\label{lem:chebyshev}
A nonzero combination of $n$ distinct real exponentials has at most $n-1$
real zeros. For increasing exponents $\alpha_1<\cdots<\alpha_n$, every
square confluent evaluation determinant is positive when rows are ordered
by increasing real node and then by increasing derivative order, and columns
by increasing exponent. For any root multiset of total multiplicity $n-1$,
the interpolants in these $n$ exponentials form a one-dimensional space;
their coefficients are nonzero and alternate in sign. They have precisely
the prescribed real roots and multiplicities.
\end{lemma}
\begin{proof}
Factor out the smallest exponential and differentiate. Rolle's theorem,
including the $m-1$ retained zeros at a root of multiplicity $m$, reduces the
bound to $n-1$ exponentials. Induction bounds any finite set of roots, hence
all real roots. Linear independence follows also from the Wronskian
\[
 e^{(\alpha_1+\cdots+\alpha_n)t}
             \prod_{i<j}(\alpha_j-\alpha_i)>0.
\]
The zero bound makes every square confluent evaluation matrix nonsingular.
For distinct ordered nodes the determinant has constant sign on their
connected ordering region. As all nodes tend to $t$, Taylor expansion gives
\[
 \frac{\det[e^{\alpha_i\tau_j}]_{i,j=1}^n}
      {\prod_{i<j}(\tau_j-\tau_i)}
 \longrightarrow
 \frac{e^{t\sum_i\alpha_i}\prod_{i<j}(\alpha_j-\alpha_i)}
      {\prod_{k=0}^{n-1}k!}>0.
\]
Partial coalescence replaces each block by successive derivative rows after
division by its positive local Vandermonde and factorials. Its limit is
nonzero by the same zero bound, and is therefore positive. For $n-1$ root
conditions, all square cofactors are positive. The null vector thus has
nonzero alternating coefficients. The global zero bound excludes further
roots and higher multiplicities.
\end{proof}

\begin{lemma}\label{lem:inward}
Let $q$ be an interpolant with $n$ distinct real exponents and $n-1$
prescribed real roots. Fix one exponent $\alpha_k$ and its coefficient.
Move another exponent $\alpha_j$ strictly toward $\alpha_k$ without crossing
an exponent, and let $P$ have the same roots and the same coefficient at
$\alpha_k$. Then
\begin{equation}\label{eq:inward}
 |P(t)|\le|q(t)|\qquad(t\in\R).
\end{equation}
\end{lemma}
\begin{proof}
The alternating coefficients in Lemma~\ref{lem:chebyshev}, together with the
shared coefficient, give $P$ and $q$ the same sign near $+\infty$. Their
exact root multiplicities agree, so their signs agree on every complementary
interval.

The difference $d=P-q$ cancels the shared term. It uses $n$ distinct
exponents: the $n-2$ unchanged ones and the old and new moved exponents.
The old exponent still has index $j$ in this ordered list. If $j<k$, insertion
and deletion occur above it; if $j>k$, deletion below and insertion just below
it cancel their changes of index. Its coefficient in $d$ is the negative of
its coefficient in $q$. All $n-1$ prescribed zeros remain zeros of $d$.
Coefficient alternation therefore gives $d$ the opposite sign to $q$ near
$+\infty$, and then on each complementary interval. Thus $P$ lies between
zero and $q$ there. All three vanish at the roots, proving \eqref{eq:inward}.
\end{proof}

The magnitude bound alone is insufficient if the coefficient of the new
exponential tends to zero. A normalized determinant estimate will prevent
that loss. Write $V(z)=\prod_{i<j}(z_j-z_i)$ and set
\[
 \mathcal H_n(\alpha,\tau)
   =\frac{\det[e^{\alpha_i\tau_j}]_{i,j=1}^n}{V(\alpha)V(\tau)}.
\]
It is positive for distinct real lists and symmetric in each list. At repeated
nodes it has a continuous positive extension given by the confluent
determinants above. Empty Vandermondes have value one.

\begin{lemma}\label{lem:detbound}
For nodes in $[0,T]$, increasing one exponent by $d>0$, while preserving
distinctness, gives
\begin{equation}\label{eq:detbound}
 1\le\frac{\mathcal H_n(\alpha+d e_i,\tau)}
              {\mathcal H_n(\alpha,\tau)}\le e^{dT}.
\end{equation}
The estimate also holds for repeated nodes.
\end{lemma}
\begin{proof}
For distinct ordered nodes a determinant integration gives the recursion
\begin{align}\label{eq:detrecursion}
 \mathcal H_n(\alpha,\tau)
 &=\frac1{V(\tau)}
 \int_{\tau_1}^{\tau_2}\cdots\int_{\tau_{n-1}}^{\tau_n}
 V(u)\mathcal H_{n-1}(\alpha_1,\ldots,\alpha_{n-1};u)
 e^{\alpha_n(\sum_j\tau_j-\sum_i u_i)}\,du .
\end{align}
Here $\mathcal H_1(\alpha_1,\tau_1)=e^{\alpha_1\tau_1}$, and any chosen
exponent can occupy the last position. To verify the recursion, put
$\beta_i=\alpha_i-\alpha_n$. Integrating the determinant
$\det[e^{\beta_i u_j}]_{i,j<n}$ column by column gives the entries
$(e^{\beta_i\tau_{j+1}}-e^{\beta_i\tau_j})/\beta_i$.
In the original determinant, factor $e^{\alpha_n\tau_j}$ from column $j$,
subtract consecutive columns and expand the last row $(1,0,\ldots,0)$.
The sign $(-1)^{n+1}$ cancels the sign in
\[
 \prod_{i<n}\beta_i=(-1)^{n-1}\prod_{i<n}(\alpha_n-\alpha_i).
\]
Division by the two Vandermondes yields \eqref{eq:detrecursion}, with no
additional factorial. Signed Vandermondes make this calculation valid for
any chosen last exponent.

On the integration region,
\[
 0\le\tau_1\le\sum_j\tau_j-\sum_i u_i\le\tau_n\le T.
\]
By symmetry, place the exponent to be increased last. All other factors in
\eqref{eq:detrecursion} remain positive and unchanged. Increasing it by $d$
multiplies the integrand by a factor between $1$ and $e^{dT}$, proving the
bound. For $n=1$ the same assertion is immediate. Limits from distinct nodes
within $[0,T]$ prove the confluent version; its denominator remains positive.
\end{proof}

\section{Extremal tests and uniqueness}\label{sec:uniqueness}

\subsection{Finite moment cones at the critical time}

For $N\ge2$, let $I_N$ be the set of real parameters $s$ for which some
$w_N\in\Mplus([0,T])$ satisfies
\begin{equation}\label{eq:finite}
 \int e^{\lambda_jt}\,dw_N(t)=m_j(s):=
           \frac{e^{\lambda_jT}-s}{\lambda_j},\qquad1\le j\le N.
\end{equation}
These auxiliary parameters may be negative; the original ratio remains
$r>0$. The cone of positive moment vectors is closed because its first
coordinate bounds mass, after which weak-star compactness applies. It is
full dimensional by Lemma~\ref{lem:chebyshev}. The vector $m(1)$, represented
by $dt$, is interior: the Gram matrix
\[
 G_{ij}=\int_0^T e^{(\lambda_i+\lambda_j)t}\,dt
\]
is positive definite, so small perturbations of $m(1)$ have positive density
$1+\sum_i d_i e^{\lambda_i t}$. The first two moments give
\begin{equation}\label{eq:ratiobounds}
 \frac{9e^{pT}-e^{9pT}}8\le s\le e^{pT}.
\end{equation}
For the lower bound, integrate $e^{9pt}-e^{pt}\ge0$ against $w_N$ and use
\eqref{eq:finite}; the upper bound follows from the first moment. Consequently
$I_N=[s_N^-,s_N^+]$ is a bounded closed interval containing $r$, with
$s_N^-<1<s_N^+$. These intervals are nested.

\begin{lemma}\label{lem:dual}
There are nonzero functions
$q_N(t)=\sum_{j=1}^N c_{j,N}e^{\lambda_jt}$, nonnegative on $[0,T]$,
and parameters
\[
 s_N=s_N^-\quad(r<1),\qquad s_N=s_N^+\quad(r>1),
\]
such that $s_N\to r$. They can be normalized so that
\begin{align}
 J_N:=\int_{-\infty}^0q_N(t)\,dt
   &=\begin{cases}-1,&r<1,\\1,&r>1,\end{cases}\label{eq:normalization}\\
 \int_0^Tq_N(t)\,dt&=|1-s_N|\le C_T,\label{eq:lone}\\
 \int_{[0,T]}q_N(t)\,dw(t)&=\delta_N\longrightarrow0
       \quad\text{for every critical scalar measure }w,\label{eq:cost}
\end{align}
where $\delta_N=r-s_N$ for $r<1$, and $\delta_N=s_N-r$ for $r>1$.
\end{lemma}
\begin{proof}
The endpoint measures of \eqref{eq:finite} have uniformly bounded mass by
\eqref{eq:ratiobounds}. Their weak-star subsequential limits represent every
odd moment at the limiting endpoint parameters $s^-\le r\le s^+$.
If $r<1$ and $s^-<r$, mix the representation at $s^-$ with $dt$ using
$\theta=(1-r)/(1-s^-)\in(0,1)$:
\[
 \theta w^-+(1-\theta)dt.
\]
It represents $r$ and has positive interval density, contradicting
Lemma~\ref{lem:slack}. This argument remains valid when $s^-<0$.
For $r>1$ and $s^+>r$, the same contradiction uses
$\theta=(r-1)/(s^+-1)$. Thus the chosen endpoints tend to $r$.

The moment vector $m(s_N)$ lies on the boundary of the closed full-dimensional
cone. A supporting hyperplane gives a nonzero coefficient vector $c_N$ with
$q_N\ge0$ on $[0,T]$ and $\int q_N\,dw_N=0$ for an endpoint measure $w_N$.
For completeness, take points $z_\ell$ outside the cone tending to this
boundary point, and let $p_\ell$ be their nearest cone points. Minimizing
squared distance gives
\[
 \left\langle\frac{p_\ell-z_\ell}{\norm{p_\ell-z_\ell}},
                  c-p_\ell\right\rangle\ge0
       \quad(c\text{ in the cone}).
\]
Here $p_\ell\to m(s_N)$. A convergent subsequence of unit normals gives
a nonzero normal nonnegative on the cone and zero at $m(s_N)$, by testing
zero and $2m(s_N)$. Testing the exponential evaluation vectors gives the
asserted $q_N$.

Since every $\lambda_j>0$, its negative-time integral is
$J_N=\sum_jc_{j,N}/\lambda_j$. The endpoint moments imply
\begin{equation}\label{eq:functional}
 0=\int q_N\,dw_N
   =\int_0^Tq_N(t)\,dt+(1-s_N)J_N.
\end{equation}
The first integral is positive because $q_N$ is nonzero and nonnegative on
a nontrivial interval. Hence $J_N\ne0$ has the stated sign, and positive
rescaling gives \eqref{eq:normalization} and \eqref{eq:lone}. Substituting
the full moments of any critical $w$ into the same functional gives
\eqref{eq:cost}.
\end{proof}

\begin{lemma}\label{lem:roots}
Each $q_N$ in Lemma~\ref{lem:dual} has exactly $N-1$ real zeros counting
multiplicity. All lie in $[0,T]$. Its coefficients are nonzero.
\end{lemma}
\begin{proof}
The endpoint measure $w_N$ is nonzero, since otherwise its first two moments
would force $s_N=e^{pT}=e^{9pT}$. Its support lies in the finite zero set of
$q_N$. Let $n$ be the number of its atoms, and let $i_0,i_T\in\{0,1\}$
indicate endpoint atoms. Nonnegativity of $q_N$ requires at least
\[
 M=2n-i_0-i_T\le N-1
\]
zeros, using double interior roots and at least simple endpoint roots.

Suppose $r<1$. If $M+i_0\le N-1$, Lemma~\ref{lem:chebyshev} constructs
an interpolant $P$ in the first $M+i_0+1$ odd exponentials with double roots
at all atoms except a simple root at $T$, when present. Choose its sign
positive on $(-\infty,T)$. Both integrals in
\[
 \int_0^TP(t)\,dt+(1-s_N)\int_{-\infty}^0P(t)\,dt
\]
are positive, whereas the finite moments identify this expression with
$\int P\,dw_N=0$. Therefore $M+i_0>N-1$, forcing $i_0=1$ and $M=N-1$.

Suppose $r>1$. If $M+1-i_0\le N-1$, construct an interpolant with a simple
root at zero, double interior roots and a simple root at $T$, when present.
It uses the first $M+2-i_0$ odd exponentials. Choose it positive on $(0,T)$;
it is negative on $(-\infty,0)$. Since $1-s_N<0$, the same expression is
strictly positive and again contradicts $\int P\,dw_N=0$. Thus $i_0=0$
and $M=N-1$.

In both cases the global zero bound is exhausted by these roots. Their
multiplicities are exact, there are no additional real roots, and the
cofactor conclusion of Lemma~\ref{lem:chebyshev} gives all $N$ nonzero
coefficients.
\end{proof}

\subsection{A coefficient that survives the limit}

The nonnegative tests have bounded integrals on the physical interval, while
their negative-time integrals remain equal to $1$ or $-1$. We use these two
facts to find one fixed coefficient that does not vanish as $N$ increases.

\begin{lemma}\label{lem:coefficient}
Along a subsequence of the tests $q_N$, there are a fixed index $m$ and
$c>0$ such that $|c_{m,N}|\ge c$ eventually.
\end{lemma}
\begin{proof}
Put $E=e^{pT}>1$ and $x=e^{pt}$. Write
\[
 q_N(t)=x f_N(x),\qquad f_N(x)=\sum_{j=1}^N c_{j,N}x^{\beta_j},
 \qquad\beta_j=(2j-1)^2-1=0,8,24,\ldots.
\]
The positive exponents have summable reciprocals and gaps at least $8$.
On $[1,E]$, $f_N\ge0$ and
\begin{equation}\label{eq:powerlone}
 \int_1^Ef_N(x)\,dx=p\int_0^Tq_N(t)\,dt\le pC_T.
\end{equation}

We use three results of Borwein--Erd\'elyi \cite{BE1997}. For
$0=\gamma_0<\gamma_1<\cdots$ with $\sum_{i\ge1}\gamma_i^{-1}<\infty$,
their Theorem 5.6 bounds the supremum of a finite polynomial in these
powers on $[0,\rho]$ by its $L^q$ norm on any measurable
$A\subset[\rho,1]$ of fixed positive measure. Their Theorem 3.2, when
$\gamma_1\ge1$, bounds the derivative on $[0,1-\eta]$ by the supremum
on $[0,1]$, with a constant independent of the number of terms.
Their Theorem 6.1, under a positive lower gap bound, gives a convergent
expansion in these powers inside $[0,R)$ for every uniform limit on $[0,R]$.

For each $R<E$, choose $\max(1,R)<R'<E$. Apply the first result with
$q=1$ to $f_N(Ez)$, on $A=[R'/E,1]$. Equation \eqref{eq:powerlone}
then bounds $\norm{f_N}_{[0,R']}$ uniformly in $N$.
For every fixed $k$, the derivative $f_N^{(k)}$ belongs to the power space
with exponents $0$ and $\{\beta_j-k:\beta_j>k\}$. The positive exponents
are integers at least one, and their reciprocals are summable: for
$\beta_j\ge2k$, $(\beta_j-k)^{-1}\le2/\beta_j$, and only finitely many
terms precede this tail. Terms with $\beta_j=k$ are constant, and terms
with $\beta_j<k$ vanish. Successive derivative estimates on nested radii
between $R$ and $E$ therefore bound every fixed derivative on $[0,R]$.

The bound for the next derivative gives equicontinuity. Arzel\`a--Ascoli
and a diagonal subsequence over derivative orders and an exhaustion of
$[0,E)$ yield locally uniform limits for every derivative. Passing through
the fundamental theorem of calculus identifies them as the derivatives of
one function $f$. The closure theorem gives
$f(x)=\sum_jc_jx^{\beta_j}$ on every $[0,R)$, $R<E$. Since the exponents
are integers, this is analytic at zero: convergence at a larger positive
$y$ bounds $|c_j|y^{\beta_j}$, giving absolute uniform convergence on
smaller complex disks. Taylor coefficients are consequently unique, and
\[
 c_{j,N}=\frac{f_N^{(\beta_j)}(0)}{\beta_j!}\longrightarrow c_j
       \quad\text{for each fixed }j.
\]
Because $E>1$, convergence is uniform on $[0,1]$. The change of variables
in the negative-time integral gives
\[
 \frac1p\int_0^1f_N(x)\,dx=J_N=\pm1.
\]
Thus $f\ne0$, and some $c_m\ne0$. Along this single subsequence
$|c_{m,N}|\ge c>0$ eventually. The subsequence and $m$ depend only on the
finite tests at the fixed $r,T$.
\end{proof}

\subsection{Inserted exponents and the two rigidity statements}

Fix the subsequence and index $m$ from Lemma~\ref{lem:coefficient}.
Recall that every $q_N$ has $N-1$ roots and that its cost
$\int q_N\,dw=\delta_N$ tends to zero for every critical $w$.
For any fixed
\begin{equation}\label{eq:muinterval}
 \lambda_m<\mu<\nu_m
\end{equation}
and $N>m$, replace $\lambda_m$ in $q_N$ by $\mu$, retain all its roots
with their multiplicities, and preserve the coefficient of $\lambda_N$.
Denote the new interpolant by $P_N$ and its $\mu$ coefficient by $d_N$.
The moved exponent stays below $\lambda_{m+1}$, and moves toward the
preserved highest exponent. Lemma~\ref{lem:inward} gives
\begin{equation}\label{eq:insertbound}
 |P_N(t)|\le q_N(t)\qquad(0\le t\le T).
\end{equation}

\begin{lemma}\label{lem:insertcoef}
For every fixed $\mu$ in \eqref{eq:muinterval}, $|d_N|\ge c_\mu>0$
eventually along the chosen subsequence.
\end{lemma}
\begin{proof}
Let $D_j$ be the positive minor of the $(N-1)$-by-$N$ confluent root
matrix with column $j$ omitted. An interpolant with last coefficient $c_N$
has coefficients $c_j=c_N(-1)^{N-j}D_j/D_N$. Replacing column $m$
leaves $D_m$ unchanged. Since the last coefficient is preserved,
\begin{equation}\label{eq:cofactorratio}
 \frac{d_N}{c_{m,N}}=\frac{D_N^{\rm old}}{D_N^{\rm new}}>0.
\end{equation}
These are square minors omitting the highest exponent. Their node factors
in the confluent normalization agree and cancel. Their exponent
Vandermondes satisfy
\[
 \frac{V_{\rm new}}{V_{\rm old}}
 =\prod_{j<m}\frac{\mu-\lambda_j}{\lambda_m-\lambda_j}
  \prod_{m<j\le N-1}\frac{\lambda_j-\mu}{\lambda_j-\lambda_m}
 \le A_m(\mu):=\prod_{j<m}\frac{\mu-\lambda_j}{\lambda_m-\lambda_j}.
\]
All factors are positive, and every factor in the second product is less
than one. The normalized determinant ratio is at most
$e^{(\mu-\lambda_m)T}$ by Lemma~\ref{lem:detbound}. Therefore
\begin{equation}\label{eq:dnonzero}
 |d_N|\ge\frac{e^{-(\mu-\lambda_m)T}}{A_m(\mu)}|c_{m,N}|
       \ge c_\mu>0.
\end{equation}
The empty product for $m=1$ is one. The constants may depend on the fixed
$\mu,m,T$, but not on $N$.
\end{proof}

We will use the following elementary uniqueness fact. If a finite signed
measure $\eta$ on $[0,T]$ satisfies
$\int e^{zt}\,d\eta(t)=0$ for all real $z$ in a nonempty open interval,
then $\eta=0$. At an interior point $z_0$, termwise differentiation is
justified by the exponential series and the finite-variation majorant
$e^{|z_0|T}e^{|z-z_0|T}\norm{\eta}_{\rm TV}$. All derivatives vanish,
so $\int t^ke^{z_0t}\,d\eta=0$ for every $k\ge0$. Uniform approximation
by real polynomials gives zero integral against every continuous function
for the weighted measure $e^{z_0t}\eta$. Uniqueness of a finite Radon
measure from continuous tests, followed by multiplication by
$e^{-z_0t}$, gives the assertion.

\begin{proposition}\label{prop:rigidity}
There is exactly one critical scalar measure. Moreover, if $w$ is critical,
$|h|\le w$, and $\int e^{\nu_jt}\,dh(t)=0$ for every $j\ge1$, then $h=0$.
\end{proposition}
\begin{proof}
For two critical measures let
$v=(w_1+w_2)/2$ and $\sigma=(w_1-w_2)/2$. Then $v$ is critical,
$|\sigma|\le v$, and all odd exponential moments of $\sigma$ vanish.
Every term of $P_N$ except its $\mu$ term belongs to that odd family.
Equations \eqref{eq:cost}, \eqref{eq:insertbound} and \eqref{eq:dnonzero}
give
\[
 |d_N|\abs{\int e^{\mu t}\,d\sigma(t)}
 =\abs{\int P_N(t)\,d\sigma(t)}
 \le\int q_N(t)\,dv(t)=\delta_N\longrightarrow0.
\]
Thus the exponential transform of $\sigma$ vanishes at every fixed $\mu$
in $(\lambda_m,\nu_m)$, and the uniqueness fact gives $\sigma=0$.

For the even family, keep the coefficient $d_N$ at $\mu$ and move the
remaining exponents toward it:
\[
 \lambda_j\mapsto\nu_j\quad(j<m),\qquad
 \lambda_j\mapsto\nu_{j-1}\quad(j>m).
\]
Perform the lower moves in increasing index order and the upper moves in
decreasing index order. The interlacing
$\lambda_j<\nu_j<\lambda_{j+1}$ ensures that no exponent is crossed.
At the central gap, $\nu_{m-1}<\mu<\nu_m$ (with the lower inequality
omitted for $m=1$). Repeated use of Lemma~\ref{lem:inward} produces an
interpolant $Q_N$ with exponents $\mu,\nu_1,\ldots,\nu_{N-1}$, the same
coefficient $d_N$ at $\mu$, and
$|Q_N|\le|P_N|\le q_N$ on $[0,T]$. For the signed measure in the
statement,
\[
 |d_N|\abs{\int e^{\mu t}\,dh(t)}
 =\abs{\int Q_N(t)\,dh(t)}
 \le\int q_N(t)\,dw(t)=\delta_N\longrightarrow0.
\]
Its transform also vanishes on $(\lambda_m,\nu_m)$, hence $h=0$.
The subsequence was fixed before choosing either measure or $\mu$; each
fixed $\mu$ uses its own positive lower bound independent of $N$.
\end{proof}

\begin{proof}[Proof of Theorem~\ref{thm:main}]
Lemma~\ref{lem:attainment} gives an admissible pair at the original bounded
minimum time. The parity map \eqref{eq:parity} sends every such pair to a
critical scalar measure $w$ and a signed $h$ dominated by $w$, with every
even moment zero. Proposition~\ref{prop:rigidity} makes $w$ unique and
forces $h=0$. Its inverse is the unique pair $(bw,bw)$. All measure
integrals and test inequalities were on the closed interval, so the conclusion
includes initial and terminal atoms.
\end{proof}

\paragraph{Use of AI tools.}
OpenAI Codex assisted with mathematical exploration, proof drafting, source
checking and editorial revision. This disclosure is not a claim of independent
human verification; the argument is provided in full for mathematical scrutiny.

\begin{thebibliography}{9}
\bibitem{LTZ2017}
J.~Loh\'eac, E.~Tr\'elat and E.~Zuazua,
\emph{Minimal controllability time for the heat equation under unilateral
state or control constraints},
Math. Models Methods Appl. Sci. \textbf{27} (2017), 1587--1644.
\href{https://doi.org/10.1142/S0218202517500270}{doi:10.1142/S0218202517500270}.

\bibitem{L2021Heat}
J.~Loh\'eac,
\emph{Nonnegative boundary control of 1D linear heat equations},
Vietnam J. Math. \textbf{49} (2021), 845--870.
\href{https://doi.org/10.1007/s10013-021-00497-5}{doi:10.1007/s10013-021-00497-5}.

\bibitem{LTZ2021}
J.~Loh\'eac, E.~Tr\'elat and E.~Zuazua,
\emph{Nonnegative control of finite-dimensional linear systems},
Ann. Inst. H. Poincar\'e C Anal. Non Lin\'eaire \textbf{38} (2021), 301--346.
\href{https://doi.org/10.1016/j.anihpc.2020.07.004}{doi:10.1016/j.anihpc.2020.07.004}.

\bibitem{BE1997}
P.~Borwein and T.~Erd\'elyi,
\emph{Generalizations of M\"untz's theorem via a Remez-type inequality for
M\"untz spaces},
J. Amer. Math. Soc. \textbf{10} (1997), 327--349.
\href{https://doi.org/10.1090/S0894-0347-97-00225-7}{doi:10.1090/S0894-0347-97-00225-7}.
\end{thebibliography}
\end{document}
